% GENERATED FILE. Edit canonical lessons, then run npm run generate:books.
% canonical-source-hash: fnv1a64:fc5c50363f8c5014
% canonical-lessons: MR-C399-calak, MR-C399-vahak, MR-C399-calavne, MR-C399-sahapravasi, MR-C399-vatsaru

\chapter{Drivers and Conductors}
\label{ch:mr-drivers-and-conductors}
\hlchaptermodality{399}

\begin{chapteropening}
\textbf{By the end of this chapter:} \emph{I can talk about a driver, a conductor and the people I travel with.}

Complete the last lesson of chapter 399: I can talk about a driver, a conductor and the people I travel with.
\end{chapteropening}

\section[\texorpdfstring{\mr{चालक} (cālak)}{चालक (cālak)}]{\mr{चालक} (cālak) — a driver}

\label{lesson:MR-C399-calak}



\begin{quote}
Before the new one: say the Marathi for a destination, then the Marathi for to set aside.
\end{quote}

\subsection*{You'll want to know: \mr{चालक}}
\textbf{\mr{चालक}} — \emph{cālak} — \textquotedblleft{}a driver\textquotedblright{}.

\begin{quote}
\mr{चालकाने} \mr{बस} \mr{हळूहळू} \mr{थांबवली}.
\end{quote}

\begin{quote}
\emph{cālakāne bas haḷūhaḷū thāmbavlī.}
\end{quote}

\begin{quote}
The driver brought the bus slowly to a stop.
\end{quote}

Masculine.

The agent of a past transitive verb takes -\mr{ने}: \mr{चालकाने} \mr{बस} \mr{थांबवली}.

\begin{cousinweb}
\textbf{\mr{चालक}} (\emph{cālak}) is a tatsama formation on \emph{cāl-}, the causative stem of Sanskrit \emph{cal}, \textquotedblleft{}to move\textquotedblright{}, with the agent ending \emph{-aka}: the one who makes the vehicle move. The same stem gives \emph{cālana}, \textquotedblleft{}causing to move\textquotedblright{} (Monier-Williams).
\end{cousinweb}

\begin{grammarlens}[title={What you've built}]
One more word for getting around.
\end{grammarlens}

\subsection*{Your turn}
\begin{itemize}
\raggedright
  \item \emph{Say it:} \emph{cālak}
  \item \emph{Say it:} \emph{cālak}, once more
  \item \emph{Say it:} \emph{cālakāne bas haḷūhaḷū thāmbavlī.} — The driver brought the bus slowly to a stop.
  \item \emph{Recall:} say the Marathi for a destination, then the Marathi for to set aside, then say \emph{cālak} again
\end{itemize}

\begin{tcolorbox}[breakable,colback=teal!4,colframe=teal!35!black,title={Before you move on}]
What does \mr{चालक} mean? (\textquotedblleft{}A driver\textquotedblright{}.) Say it once more.
\end{tcolorbox}
\section[\texorpdfstring{\mr{वाहक} (vāhak)}{वाहक (vāhak)}]{\mr{वाहक} (vāhak) — a bus conductor}

\label{lesson:MR-C399-vahak}



\begin{quote}
Before the new one: say the Marathi for to set aside, then the Marathi for a driver.
\end{quote}

\subsection*{You'll want to know: \mr{वाहक}}
\textbf{\mr{वाहक}} — \emph{vāhak} — \textquotedblleft{}a bus conductor\textquotedblright{}.

\begin{quote}
\mr{वाहकाने} \mr{माझं} \mr{तिकीट} \mr{तपासलं}.
\end{quote}

\begin{quote}
\emph{vāhakāne mājhaṁ tikīṭ tapāslaṁ.}
\end{quote}

\begin{quote}
The conductor checked my ticket.
\end{quote}

Masculine.

On a bus the \mr{चालक} drives and the \mr{वाहक} sells and checks tickets.

\begin{cousinweb}
\textbf{\mr{वाहक}} (\emph{vāhak}) is a tatsama, taken from Sanskrit \emph{vāhaka}, \textquotedblleft{}one who carries, a carrier\textquotedblright{} (Monier-Williams); on a bus it names the conductor.
\end{cousinweb}

\begin{grammarlens}[title={What you've built}]
One more word for getting around.
\end{grammarlens}

\subsection*{Your turn}
\begin{itemize}
\raggedright
  \item \emph{Say it:} \emph{vāhak}
  \item \emph{Say it:} \emph{vāhak}, once more
  \item \emph{Say it:} \emph{vāhakāne mājhaṁ tikīṭ tapāslaṁ.} — The conductor checked my ticket.
  \item \emph{Recall:} say the Marathi for to set aside, then the Marathi for a driver, then say \emph{vāhak} again
\end{itemize}

\begin{tcolorbox}[breakable,colback=teal!4,colframe=teal!35!black,title={Before you move on}]
What does \mr{वाहक} mean? (\textquotedblleft{}A bus conductor\textquotedblright{}.) Say it once more.
\end{tcolorbox}
\section[\texorpdfstring{\mr{चालवणे} (cālavṇe)}{चालवणे (cālavṇe)}]{\mr{चालवणे} (cālavṇe) — to drive; to run}

\label{lesson:MR-C399-calavne}



\begin{quote}
Before the new one: say the Marathi for a driver, then the Marathi for a bus conductor.
\end{quote}

\subsection*{You'll want to know: \mr{चालवणे}}
\textbf{\mr{चालवणे}} — \emph{cālavṇe} — \textquotedblleft{}to drive; to run\textquotedblright{}.

\begin{quote}
\mr{माझे} \mr{काका} \mr{बस} \mr{चालवतात}.
\end{quote}

\begin{quote}
\emph{mājhe kākā bas cālavtāt.}
\end{quote}

\begin{quote}
My uncle drives a bus.
\end{quote}

Past, with -\mr{ने} (he …): \mr{त्याने} \mr{चालवलं}, \emph{tyāne cālavlaṁ}.

Not the same as \mr{चालणे}, to walk: \mr{चालवणे} is to make something go, so \mr{गाडी} \mr{चालवणे} is to drive and \mr{दुकान} \mr{चालवणे} is to run a shop.

\begin{cousinweb}
\textbf{\mr{चालवणे}} (\emph{cālavṇe}) is built on \textbf{\mr{चालणे}}, \textquotedblleft{}to walk, to go\textquotedblright{}, which CDIAL lists as Marathi \emph{ċālṇẽ} under \emph{*calyati}, \textquotedblleft{}goes\textquotedblright{} (CDIAL 4721), as its causative: to make a vehicle or a business go.
\end{cousinweb}

\begin{grammarlens}[title={What you've built}]
One more word for getting around.
\end{grammarlens}

\subsection*{Your turn}
\begin{itemize}
\raggedright
  \item \emph{Say it:} \emph{cālavṇe}
  \item \emph{Say it:} \emph{cālavṇe}, once more
  \item \emph{Say it:} \emph{mājhe kākā bas cālavtāt.} — My uncle drives a bus.
  \item \emph{Recall:} say the Marathi for a driver, then the Marathi for a bus conductor, then say \emph{cālavṇe} again
\end{itemize}

\begin{tcolorbox}[breakable,colback=teal!4,colframe=teal!35!black,title={Before you move on}]
What does \mr{चालवणे} mean? (\textquotedblleft{}To drive; to run\textquotedblright{}.) Say it once more.
\end{tcolorbox}
\section[\texorpdfstring{\mr{सहप्रवासी} (sahapravāsī)}{सहप्रवासी (sahapravāsī)}]{\mr{सहप्रवासी} (sahapravāsī) — a fellow traveller}

\label{lesson:MR-C399-sahapravasi}



\begin{quote}
Before the new one: say the Marathi for a bus conductor, then the Marathi for to drive.
\end{quote}

\subsection*{You'll want to know: \mr{सहप्रवासी}}
\textbf{\mr{सहप्रवासी}} — \emph{sahapravāsī} — \textquotedblleft{}a fellow traveller\textquotedblright{}.

\begin{quote}
\mr{माझ्या} \mr{सहप्रवाशाने} \mr{मला} \mr{पाणी} \mr{दिलं}.
\end{quote}

\begin{quote}
\emph{mājhyā sahapravāśāne malā pāṇī dilaṁ.}
\end{quote}

\begin{quote}
My fellow traveller gave me some water.
\end{quote}

Masculine.

\mr{सह}- means together with; before -\mr{ने} and -\mr{ला} it becomes \mr{सहप्रवाशा}-.

\begin{cousinweb}
\textbf{\mr{सहप्रवासी}} (\emph{sahapravāsī}) is a compound of Sanskrit \emph{saha}, \textquotedblleft{}together with\textquotedblright{} (Monier-Williams), and \textbf{\mr{प्रवासी}}, \textquotedblleft{}a traveller\textquotedblright{}, built on Sanskrit \emph{pravāsa}, \textquotedblleft{}a journey abroad\textquotedblright{}.
\end{cousinweb}

\begin{grammarlens}[title={What you've built}]
One more word for getting around.
\end{grammarlens}

\subsection*{Your turn}
\begin{itemize}
\raggedright
  \item \emph{Say it:} \emph{sahapravāsī}
  \item \emph{Say it:} \emph{sahapravāsī}, once more
  \item \emph{Say it:} \emph{mājhyā sahapravāśāne malā pāṇī dilaṁ.} — My fellow traveller gave me some water.
  \item \emph{Recall:} say the Marathi for a bus conductor, then the Marathi for to drive, then say \emph{sahapravāsī} again
\end{itemize}

\begin{tcolorbox}[breakable,colback=teal!4,colframe=teal!35!black,title={Before you move on}]
What does \mr{सहप्रवासी} mean? (\textquotedblleft{}A fellow traveller\textquotedblright{}.) Say it once more.
\end{tcolorbox}
\section[\texorpdfstring{\mr{वाटसरू} (vāṭsarū)}{वाटसरू (vāṭsarū)}]{\mr{वाटसरू} (vāṭsarū) — a wayfarer; a passer-by}

\label{lesson:MR-C399-vatsaru}



\begin{quote}
Before the new one: say the Marathi for to drive, then the Marathi for a fellow traveller.
\end{quote}

\subsection*{You'll want to know: \mr{वाटसरू}}
\textbf{\mr{वाटसरू}} — \emph{vāṭsarū} — \textquotedblleft{}a wayfarer; a passer-by\textquotedblright{}.

\begin{quote}
\mr{एका} \mr{वाटसरूने} \mr{आम्हाला} \mr{रस्ता} \mr{सांगितला}.
\end{quote}

\begin{quote}
\emph{ekā vāṭsarūne āmhālā rastā sāṅgitlā.}
\end{quote}

\begin{quote}
A passer-by told us the way.
\end{quote}

Masculine.

An old, warm word for someone on the road, often walking.

\begin{cousinweb}
\textbf{\mr{वाटसरू}} (\emph{vāṭsarū}) is built on \textbf{\mr{वाट}}, \textquotedblleft{}a path\textquotedblright{}, inherited from Sanskrit \emph{vartman}, \textquotedblleft{}track of a wheel, path\textquotedblright{} (CDIAL 11366); the ending marks the person on it.
\end{cousinweb}

\begin{grammarlens}[title={What you've built}]
One more word for getting around.
\end{grammarlens}

\subsection*{Your turn}
\begin{itemize}
\raggedright
  \item \emph{Say it:} \emph{vāṭsarū}
  \item \emph{Say it:} \emph{vāṭsarū}, once more
  \item \emph{Say it:} \emph{ekā vāṭsarūne āmhālā rastā sāṅgitlā.} — A passer-by told us the way.
  \item \emph{Recall:} say the Marathi for to drive, then the Marathi for a fellow traveller, then say \emph{vāṭsarū} again
\end{itemize}

\begin{tcolorbox}[breakable,colback=teal!4,colframe=teal!35!black,title={Before you move on}]
What does \mr{वाटसरू} mean? (\textquotedblleft{}A wayfarer; a passer-by\textquotedblright{}.) Say it once more.
\end{tcolorbox}
